\section{Metodolog\'{i}a}
\label{sec:methodology}

\subsection{Dise\~{n}o de la Investigaci\'{o}n}
\label{subsec:design}

[Tipo de investigaci\'{o}n: aplicada, cuantitativa, experimental. Enfoque: comparaci\'{o}n de algoritmos ML + componente prescriptivo.]

\subsection{Arquitectura del Sistema Propuesto}
\label{subsec:architecture}

[Descripci\'{o}n del pipeline:
\begin{enumerate}
    \item Ingesta y preprocesamiento de datos
    \item Feature engineering
    \item Entrenamiento y evaluaci\'{o}n de modelos predictivos
    \item M\'{o}dulo prescriptivo (optimizaci\'{o}n de inventario)
    \item Generaci\'{o}n de recomendaciones
\end{enumerate}]

\subsection{Algoritmos de Machine Learning}
\label{subsec:algorithms}

[Descripci\'{o}n formal de cada algoritmo evaluado:]

\subsubsection{Random Forest}
[Formulaci\'{o}n, hiperpar\'{a}metros clave, justificaci\'{o}n de inclusi\'{o}n.]

\subsubsection{XGBoost}
[Formulaci\'{o}n basada en \cite{chen2016xgboost}, ventajas para datos tabulares.]

\subsubsection{LightGBM}
[Formulaci\'{o}n basada en \cite{ke2017lightgbm}, eficiencia computacional.]

\subsubsection{Support Vector Regression}
[Formulaci\'{o}n, kernel selection.]

\subsubsection{Baselines Estad\'{i}sticos}
[ARIMA, Holt-Winters como benchmarks de comparaci\'{o}n.]

\subsection{M\'{e}tricas de Evaluaci\'{o}n}
\label{subsec:metrics}

Las m\'{e}tricas utilizadas para evaluar el desempe\~{n}o predictivo son:

\begin{align}
    \text{RMSE} &= \sqrt{\frac{1}{n}\sum_{i=1}^{n}(y_i - \hat{y}_i)^2} \label{eq:rmse} \\
    \text{MAE}  &= \frac{1}{n}\sum_{i=1}^{n}|y_i - \hat{y}_i| \label{eq:mae} \\
    \text{MAPE} &= \frac{100}{n}\sum_{i=1}^{n}\left|\frac{y_i - \hat{y}_i}{y_i}\right| \label{eq:mape}
\end{align}

[Para el componente prescriptivo: costo total de inventario, nivel de servicio, fill rate.]

\subsection{Estrategia de Validaci\'{o}n}
\label{subsec:validation}

[Walk-forward validation (expanding window) para respetar la estructura temporal. Justificaci\'{o}n de por qu\'{e} no usar random split.]
